\chapter{考察}

\section{本章の概要}

本章では、実験1から実験3を踏まえた総合的な考察を述べる。

\section{結果の解釈}

提案手法が分岐分類精度と Branch Contrast Accuracy を改善する傾向を示した場合、その意義は二つある。第一に、VLM に単に「次にどうするか」を答えさせるよりも、不確実性を証拠充足性、不確実性原因、工程関連性に分解して診断させることが有効であると解釈できる。第二に、同じ視覚状態でも工程文脈に応じて適切な分岐が変わるという、長期タスク固有の問題設定が妥当であることを裏づけられる。

\section{手法の限界}

本研究は、机上片付けを想定した限定的な長期タスクに基づくため、対象物、環境、ロボット身体が変わった場合の汎化性は十分に検証できない。また、VLM の自然言語理由は後付け説明である可能性があるため、理由文そのものではなく、構造化ラベルと分岐結果を中心に評価する必要がある。さらに、安全性を重視するほど ask\_user や stop が増え、作業効率が低下するおそれがある。

\section{今後の方向性}

今後は、複数タスクへの拡張、VLA との接続、再観測視点の自動選択、ユーザ確認文の生成、実行後の回復行動との統合を検討する。また、world model や計画器と連携し、defer した不確実性を後工程で回収する仕組みへ拡張することも重要である。
